%====================== CONCLUSION GÉNÉRALE ======================
\chapter*{Conclusion générale}
\addcontentsline{toc}{chapter}{Conclusion générale}
\markboth{Conclusion générale}{Conclusion générale}

Ce projet de fin d'études a abouti à la conception et à la réalisation de
\devpilot{}, une plateforme agentique de RAG de qualité production. L'objectif
initial — combler l'écart entre la puissance générative des grands modèles de
langage et l'exigence de réponses \emph{fondées}, \emph{citées}, \emph{mesurées}
et \emph{traçables} — est atteint et validé sur l'ensemble du pipeline.

\paragraph{Réalisations.} La plateforme centralise les connaissances en espaces
isolés, les ingère de façon asynchrone, les indexe vectoriellement, et restitue
des réponses ancrées et citées via un flux agentique de sept agents diffusé en
\emph{streaming}. Chaque réponse est évaluée et, le cas échéant, corrigée.
L'ensemble est instrumenté par trois contributions distinctives : la tour de
contrôle \textbf{RAGOps}, le sous-système de \textbf{traçabilité} (explorateur et
inspecteur de traces) et l'\textbf{AI Execution Studio}.

\paragraph{Maturité architecturale.} L'architecture en couches, asynchrone et
neutre vis-à-vis des fournisseurs d'IA, la discipline de typage statique de bout
en bout, la conteneurisation complète et la fondation de tests confèrent au
système une maturité dépassant celle d'un prototype académique.

\paragraph{Compétences acquises.} Ce travail a consolidé des compétences en
architecture des systèmes distribués, en ingénierie des plateformes d'IA (RAG,
agents, évaluation), en observabilité et traçabilité, ainsi qu'en ingénierie
\emph{frontend} avancée et en pratiques DevOps.

\paragraph{Enseignement principal.} La valeur d'une plateforme d'IA d'entreprise
ne réside pas seulement dans la qualité de ses réponses, mais dans leur
\emph{auditabilité} et leur \emph{exploitabilité}. Concevoir pour l'observabilité
dès l'origine et raisonner systématiquement en compromis d'ingénierie ont été les
leviers déterminants de ce projet.

%====================== TRAVAUX FUTURS ======================
\chapter*{Perspectives et travaux futurs}
\addcontentsline{toc}{chapter}{Perspectives et travaux futurs}
\markboth{Perspectives}{Perspectives}

Plusieurs axes d'évolution prolongent naturellement ce travail.

\begin{description}[leftmargin=2.4cm,style=nextline]
\item[Orchestration multi-agents] Adoption d'un graphe d'exécution explicite
(type LangGraph) pour modéliser les transitions, les boucles et les stratégies de
repli du flux agentique.
\item[RAG avancé] Récupération hybride enrichie, fenêtre de contexte adaptative,
et \textbf{GraphRAG} exploitant un graphe de connaissances pour les questions
multi-sauts.
\item[Évaluation formelle] Intégration de RAGAS~\cite{es2023} et de jeux
d'évaluation versionnés pour un suivi longitudinal de la qualité.
\item[MCP et connecteurs d'entreprise] Intégration du \emph{Model Context
Protocol} et de connecteurs (dépôts de code, wikis, tickets) pour étendre les
sources de connaissance.
\item[Agents autonomes] Agents capables de planifier des actions multi-étapes et
d'invoquer des outils sous contrôle de gouvernance.
\item[RAGOps avancé et gouvernance de l'IA] Détection d'anomalies, budgets de
coût et de latence, politiques de rétention, et pistes d'audit étendues pour la
gouvernance de l'IA.
\item[Déploiement cloud et multi-tenant] Migration vers une orchestration
Kubernetes, cloisonnement organisationnel, quotas et facturation par locataire.
\end{description}

Ces perspectives s'appuient sur les fondations posées — traçabilité native,
architecture découplée et observabilité — et visent à faire évoluer \devpilot{}
d'une plateforme de démonstration aboutie vers un produit d'entreprise à grande
échelle.
